\chapter*{R\'{e}sum\'{e}}
\addcontentsline{toc}{chapter}{R\'{e}sum\'{e}}
\markboth{R\'{e}sum\'{e}}{}

\vspace{0.5cm}

Face à l'essor des solutions de formation immersive en entreprise, ce projet de fin
d'études, réalisé au sein de Tynass IT, porte sur la conception et le développement
de \tynassit{}, une plateforme B2B dédiée à la gestion et à la diffusion de formations
professionnelles en Réalité Virtuelle (VR) sur le casque Meta Quest~3.

La plateforme repose sur trois composantes complémentaires : un backend REST
développé avec Node.js, Express et MongoDB Atlas ; deux dashboards web développés
avec React et TypeScript ; et une application VR Unity~6 déployée sur Meta Quest~3,
permettant aux employés de vivre des simulations immersives de formation.

Dans le cadre du cas d'application pilote réalisé avec Indigo Company, un module
d'intégration des nouveaux employés a été développé : exploration des locaux,
présentation des informations internes et évaluation des connaissances sous forme
de quiz, avec envoi automatique des résultats vers la plateforme. Cette approche
permet de faciliter le suivi des résultats par l'entreprise et de structurer davantage
le parcours d'intégration des nouveaux employés.

\textbf{Mots-clés :} Réalité Virtuelle, VR, B2B, formation professionnelle,
Meta Quest~3, Unity~6, Node.js, React, MongoDB, JWT, PFE.

\clearpage
